%------------------------------------------------------------------------------%
% \chapter{Matrizes}
%------------------------------------------------------------------------------%

%------------------------------------------------------------------------------%
\begin{exercicios}{Sistemas Lineares}
    
    \begin{exercicio}
        Diga qual das equações abaixo são lineares
        \begin{mathlist}
            \item x_1-x_2+x_3-2x_4 =3
            \item x_1 + mx_2+x_3^2=n \) \\[1mm] onde $m$ e $n$ são constantes dadas
            \item x-2y+3z=4
            \item a_1x_1 + a_2x_2^2+a_3x_3^3=b \) \\[1mm] onde $a_i$ e $b_i$ são constantes dadas
            \item 2x_1 + \log(x_2) +x_3 =\log(2)
            \item -x_1 +x_2 -3x_3 -x_4 - x_5 = 0
            \item \sqrt{3}x_1 +\sqrt{2}x_2 +x_3 = 5
            \item x_1 + 3x_2 - 4x_3 + 5x_4 = 10 - 2x_5
            \item e^2x_1 + \log(4)x_2 = \sqrt{3}
            \item x_1 + x_1x_2 = x_3
        \end{mathlist}
    \end{exercicio}
    
    \begin{exercicio}
        Verifique se $(2,0,-3)$ é solução de \[ 2x_1+5x_2+2x_3=-2 \]
    \end{exercicio}
    
    \begin{exercicio}
        Verifique se $(1,1,-1,-1)$ é solução de \[ 5x_1 -10x_2 -x_3 +2x_4=0 \]
    \end{exercicio}
        
    \begin{exercicio}
        Encontre uma solução para a equação linear 
        \[ 2x_1 -x_2 -x_3 = 0 \]
        diferente de $(0,0,0)$
    \end{exercicio}
        
    \begin{exercicio}
        Escreva os seguintes sistemas na forma matricial, 
        considerando que $a$, $b$, $c$, $d$, $e$, $f$, $m$ e $n$ são constantes
        \begin{mathlist}
            \item \systeme{
                x-y+z    =2,
            	-x+2y+2z =5,
            	5x-y+5z  =1}
            \item \systeme{
            	3x-5y+4z-t = 8,
            	2x+y-2z    = -3,
            	-x-2y+z-3t = 1,
            	-5x -y+6t  = 4}
            \item \systeme{
            	ax+by+cz   = d,
            	-mz + ny   = e,
            	abx-b^2+mz = f}
            \item \systeme{
            	\sqrt{2}x-3y+2z  = 7,
            	+7y -z           = 0,
            	4x+\sqrt{3}y +2z = 5}
            \item \systeme{
            	ax-by+2z  = 1,
            	a^2-by +z = 3}
            \item \systeme{
            	x+y-z   = 3-t,
            	-x-y-2z =1-3t,
            	5x+3z   =7+t}
        \end{mathlist}
    \end{exercicio}
    
    \begin{exercicio}
        Verifique se $(0,-3,-4)$ é solução do sistema
        \[\systeme{
        	x+y-z      =1,
        	2x-y+z     = -1,
        	x + 2y + z = 2}
        \]
    \end{exercicio}

    \begin{exercicio}
        Verifique se $(1,0,-2,1)$ é solução do sistema
        \[\systeme{
        	5x+3y-2z-4t  =5,
        	2x-4y +3z -5t =-9,
        	-x+2y-5z+3t   =12}
        \]
    \end{exercicio}
        
    \begin{exercicio}
        Encontre a solução do sistema $(A+4I_3)X=0$ sabendo que
        \[ A = \left[ 
        \begin{array}{rrr}
            1 & 0 & 5 \\
            1 & 1 & 1 \\
            0 & 1 & 4 \\
        \end{array}
        \right]\]
    \end{exercicio}
    
    \begin{exercicio}
        Encontre os valores de $a$ para os quais a matriz
        \[A = \left[ 
        \begin{array}{rrr}
            1 & 1 & 0 \\
            1 & 0 & 0 \\
            1 & 2 & a \\
        \end{array}
        \right]\]
        possui inversa.
    \end{exercicio}
    
    \begin{exercicio}
        Sejam
        \[A = \left[ 
            \begin{array}{rrr}
                1 & 0 & 5 \\
                1 & 1 & 1 \\
                0 & 1 & 4 \\
            \end{array}
            \right] 
            \qquad  
        B = \left[ \begin{array}{c} 5 \\ 5 \\ 10 \\ \end{array} \right] 
        \]
        encontre a solução do sistema $(A+4I_3)X=B$.
    \end{exercicio}
    
    \begin{exercicio}
        Resolva o sistema
        \[\begin{cases}
        	\dfrac{2}{u} + \dfrac{3}{v} =  8 \\[4mm]
        	\dfrac{1}{u} - \dfrac{1}{v} = -1
        \end{cases}\]
    \end{exercicio}
    
    \begin{exercicio}
        Os valores de $x$, $y$ e $z$, solução do sistema
        \[\systeme{
        	x + 2y + 3z  = 14,
        	4x + 5y + 6z = 32,
        	7x + 8y + 9z = a}
        \]
        formam, nesta ordem, uma P.A. de razão $1$. Qual é o valor de $a$?	
    \end{exercicio}
    
    \begin{exercicio}
        Resolva o sistema utilizando a regra de Cramer
        \[\systeme{
            4x - 3y = 11,
            2x + 5y = -1}
        \]
    \end{exercicio}
    
    \begin{exercicio}
        Resolva o sistema a seguir utilizando a regra de Cramer
        \[\systeme{
        	2x + y - z  = 5,
        	3x - 2y + z = -2,
        	x + z       = 0}
        \]
    \end{exercicio}
    
    \begin{exercicio}
        Discuta o sistema, em função dos valores do parâmetro real $a$.
        \[\systeme{
            a^2x + y = 1,
            x + y    = a}
        \]
    \end{exercicio}
    
    \begin{exercicio}
        Discuta, por Cramer, o sistema
        \[\systeme{
            7x + y - 3z = 10,
            x + y + z   = 6,
            4x + y + Pz = Q}
        \]
    \end{exercicio}
    
    \begin{exercicio}
        Encontre o valor de $a$ para que o sistema seja possível. 
        \[\systeme{
            2x - y + 3z  = a,
            x + 2y - z   = 3,
            7x + 4y + 3z = 13}
        \]
        Para o valor encontrado de $a$, ache a solução geral do sistema, isto é, 
        ache expressões que representem todas as soluções do sistema. 
        Explicite duas dessas soluções.		
    \end{exercicio}
    
    \begin{exercicio}
        Calcule o valor de $m$ para que o sistema admita soluções diferentes da trivial
        \[\systeme{
            x + y - z   = 0,
            2x - y + z  = 0,
            3x + my - z = 0}
        \]
    \end{exercicio}
    
    \begin{exercicio}
        Sejam $a$, $b$ e $c$ números tais que
        \[\systeme{
            a + 2b = 5,
            b + 2c = 8,
            2a + c = 5}
        \]
        Determine o valor de $a + b + c$.		
    \end{exercicio}
    
    \begin{exercicio}
        Calcule $a$ e $b$ de modo que o sistema a seguir seja indeterminado
        \[\systeme{
            x + 2y + 2z  = a,
            3x + 6y - 4z = 4,
            2x + by - 6z = 1}
        \]
    \end{exercicio}
    
    \begin{exercicio}
        Resolva o sistema
        \[\systeme{
            2y + x = b,
            2z - y = b,
            az + x = b}
        \]
        para os casos
        \begin{mathlist}[2]
            \item a = 0 \text{ e } b = 1
            \item a = 4 \text{ e } b = 0
        \end{mathlist}		
    \end{exercicio}
    
    \begin{exercicio}
        Ache os valores de $k \in \R$ de maneira que o sistema linear seja impossível.
        \[\systeme{
             x + y = 1,
             y + z = 2,
            kz + x = 0}
        \]
    \end{exercicio}
    
\end{exercicios} % Sistemas Lineares

%------------------------------------------------------------------------------%
\begin{exercicios}{Método de Gauss}

    \begin{exercicio}
        Resolva o sistema pelo método de Gauss
        \[ \systeme{  x + 3y = 7,
                     2x +  y = 4} 
        \]
        \begin{answers}
            \(\qquad x=1 \qquad y=2 \)
        \end{answers}
    \end{exercicio}
    
    \begin{exercicio}
        Resolva o sistema pelo método de Gauss
        \[\systeme{ 2x + y -2z = 10,
            3x +2y +2z =  1,
            5x +4y +3z =  4}
        \]
        \begin{answers}
            \(\qquad x=1 \qquad y=2 \qquad z=-3 \)
        \end{answers}
    \end{exercicio}
    
    \begin{exercicio}
        Resolva o seguinte sistema de equações lineares
        \[\systeme{2x +  y +  z +  w = 1,
            x + 2y +  z +  w = 2,
            x +  y + 2z +  w = 3,
            x +  y +  z + 2w = 4}
        \]
        \begin{answers}
            \(\qquad w=-1 \qquad x=0 \qquad y=1 \qquad z=2 \)
        \end{answers}
    \end{exercicio}
    
    \begin{exercicio}
        Resolva o sistema pelo método de Gauss
        \[\systeme{ x  -y -2z = 1,
            -x  +y + z = 2,
            x -2y + z =-2}
        \]
        \begin{answers}
            \(\qquad x=-11 \qquad y=-6 \qquad z=-3 \)
        \end{answers}
    \end{exercicio}
    
    \begin{exercicio}
        Resolva o sistema pelo método de Gauss
        \[\systeme{  x -2y +3z =  1,
            3x +6y -3z = -2,
            6x +6y +3z =  5}
        \]
        \begin{answers}
            \(\qquad x=2 \qquad y=\sfrac{7}{4} \qquad z=\sfrac{13}{6} \)
        \end{answers}
    \end{exercicio}
    
    \begin{exercicio}
        Use método de Gauss para resolver o sistema linear $(A+2I)X=B$, onde
        \[A = \left[ \begin{array}{rrr}
        1 & 0 & 5 \\
        1 & 1 & 1 \\
        0 & 1 & 4
        \end{array} \right] 
        \qquad
        B = \left[ \begin{array}{r} 3 \\ 4 \\ 1 \\ \end{array}
        \right] 
        \]
        \begin{answers}
            \(
            \qquad A+2I = \left[
            \begin{array}{rrr}
            3 & 0 & 5 \\
            1 & 3 & 1 \\
            0 & 1 & 6
            \end{array}
            \right]
            \qquad X = \left[
            \begin{array}{r}
            1 \\
            1 \\
            0
            \end{array}
            \right]
            \)
        \end{answers}
    \end{exercicio}
            
    \begin{exercicio}
        Resolva se possível o sistema pelo método de Gauss
        \[\systeme{ x - 2y +  z = 1,
                   2x +  y -  z = 2,
                    x + 3y - 2z = 1}
        \]
        \begin{answers}
            \(\qquad\) Sistema possui infinitas soluções.
        \end{answers}
    \end{exercicio}
    
    \begin{exercicio}
        Resolva se possível o sistema pelo método de Gauss
        \[\systeme{ x + 3y -  z = 6,
                   2x + 4y + 2z = 9,
                    x +  y + 3z = 1}
        \]
        \begin{answers}
            \(\qquad\) Sistema não possui solução.
        \end{answers}
    \end{exercicio}
    
    \begin{exercicio}
        Resolva se possível o sistema pelo método de Gauss
        \[\systeme{ 3x + 4y -  z = 1,
                    4x + 5y + 2z = 12,
                     x - 2y + 3z = 8}
        \]
        \begin{answers}
            \(\qquad x=0,1 \qquad y=1 \qquad z=3,3 \)
        \end{answers}
    \end{exercicio}
    
    \begin{exercicio}
        Resolva se possível o sistema pelo método de Gauss
        \[\systeme{  x + 2y -3z = 1,
            3x -  y +2z = 7,
            5x + 3y -4z = 2}
        \]
        \begin{answers}
            \(\qquad x=-7 \qquad y=79 \qquad z=50 \)
        \end{answers}
    \end{exercicio}
    
    \begin{exercicio}
        Resolva se possível o sistema pelo método de Gauss
        \[\systeme{  x +2y -3z =  6,
            2x - y +4z =  2,
            4x +3y -2z = 14}
        \]
        \begin{answers}
            \(\qquad x=\sfrac{7}{3} \qquad y=\sfrac{4}{3} \qquad z=-\sfrac{1}{3} \)
        \end{answers}
    \end{exercicio}
    
    \begin{exercicio}
        Resolva se possível o sistema pelo método de Gauss
        \[\systeme{ 2x +3y +z = 0,
            x -2y -z = 1,
            x +4y +z = 2}
        \]
        \begin{answers}
            \(\qquad x=-\sfrac{1}{4} \qquad y=\sfrac{7}{4} \qquad z=-\sfrac{19}{4} \)
        \end{answers}
    \end{exercicio}
    
    \begin{exercicio}
        Resolva se possível o sistema pelo método de Gauss
        \[\systeme{-x + y - 2z      = 1,
            2x - y +      3t = 2,
            x -2y +  z - 2t = 0}
        \]
        \begin{answers}
            \(\qquad \) Sistema possui infinitas soluções
        \end{answers}
    \end{exercicio}
    
    \begin{exercicio}
        Resolva se possível o sistema pelo método de Gauss
        \[\systeme{ x +3y  + 2z = 2,
            3x +5y +4z = 4,
            5x +3y +4z = -10}
        \]
        \begin{answers}
            \(\qquad \) Sistema não possui solução
        \end{answers}
    \end{exercicio}
    
    \begin{exercicio}
        Resolva se possível o sistema pelo método de Gauss
        \[\systeme{ x + y - z + t = 1,
            3x - y -2z + t = 2,
            -x -2y +3z +2t = -1}
        \]
        \begin{answers}
            \(\qquad \) Sistema possui infinitas soluções
        \end{answers}
    \end{exercicio}
    
    \begin{exercicio}
        Resolva se possível o sistema pelo método de Gauss
        \[\systeme{ x + y + z + t = 1,
            x - y + z + t = -1,
            y - z +2t = 2,
            2x + z - t = -1 }
        \]
        \begin{answers}
            \(\qquad x=-\sfrac{1}{5} \qquad y=1 \qquad z=-\sfrac{1}{5} \qquad t=\sfrac{2}{5} \)
        \end{answers}
    \end{exercicio}
    
    \begin{exercicio}
        Resolva se possível o sistema pelo método de Gauss
        \[\systeme{ x -2y -3z = 5,
            -2x +5y +2z = 3,
            -x +3y - z = 2}
        \]
        \begin{answers}
            \(\qquad \) Sistema não possui solução
        \end{answers}
    \end{exercicio}
    
\end{exercicios} % Método de Gauss

%------------------------------------------------------------------------------%
